\rok{Новембар 2023.}{G}

Први практични тест из Алгоритама и структура података

\zadatak{1}{Брисање елемената из реда (Queue Dequeue)}
Написати методу која имплементира стандардан начин за брисање елемената из реда.

\textbf{Додатне напомене за решавање задатка:}
\begin{itemize}
    \item Алгоритам написати у оквиру закоментарисане методе \texttt{dequeue?()}.
    \item Поменута метода налази се у оквиру класе \texttt{Queue} која се налази у придруженом шаблону.
    \item Обезбедити код у \texttt{Main} методи за тестирање алгоритма.
\end{itemize}



\zadatak{2}{Да ли елементи листе стварају Капрекаров број?}
Креирати функцију која ће проверавати да ли су елементи једноструко спрегнуте листе заправо Капрекаров број. Капрекаров број је број чији квадрат када се подели на два дела и такав да је збир делова једнак изворном броју, а нити један део нема вредност 0. Алгоритам је неопходно написати са линеарном временском комплексношћу $O(n)$.

\textbf{Додатни захтеви за решавање задатка:}
\begin{itemize}
    \item Захтеван потпис методе:
    \begin{Verbatim}[fontsize=\footnotesize, numbers=left, xleftmargin=10mm]
public static bool IsNumberKaprekar ()
    \end{Verbatim}
\end{itemize}

\textbf{Пример:}
\begin{itemize}
    \item \textbf{Улаз:} $4 \rightarrow 5$ \\
          \textbf{Објашњење:} $45^2 = 2025$; $20 + 25 = 45$ \\
          \textbf{Излаз:} \texttt{true}
    \item \textbf{Улаз:} $2 \rightarrow 9 \rightarrow 7$ \\
          \textbf{Објашњење:} $297^2 = 88209$; $88 + 209 = 297$ \\
          \textbf{Излаз:} \texttt{true}
\end{itemize}



\zadatak{3}{Исограм}
Креирање функције која проверава да ли је неки број исограм. Исограм је број који не дуплира цифре. Алгоритам је неопходно написати са линеарно-логаритамском временском комплексношћу $O(n \log n)$.

\textbf{Додатни захтеви за решавање задатка:}
\begin{itemize}
    \item Алгоритам треба да врати \texttt{bool} вредност.
    \item Као помоћ може се користити алгоритам за сортирање који је дат у шаблону задатка.
    \item Цифре прослеђеног броја неопходно је сместити у низ. Уколико се користи динамички низ, користити функцију \texttt{ToArray()} како би се од њега добио статички низ.
    \item Осим динамичког низа, може се користити и структура једноструко повезане листе која је дата у оквиру задатка 2.
    \item Захтеван потпис методе:
    \begin{Verbatim}[fontsize=\footnotesize, numbers=left, xleftmargin=10mm]
public bool IsNumberIsogram(long number)
    \end{Verbatim}
\end{itemize}

\textbf{Примери:}
\begin{itemize}
    \item \textbf{Улаз 1:} 12345 \\
          \textbf{Излаз 1:} \texttt{true}
    \item \textbf{Улаз 2:} 12324 \\
          \textbf{Излаз 2:} \texttt{false}
\end{itemize}

\sablon{Шаблон првог теста новембар 2023. група G}{2023/Novembar/GrupaG/AISP2023_Test1_GrupaG.linq}
